\begin{examplebox}
	\textbf{\textcolor{CExample}{例 1（基础）}}
	\textbf{\textcolor{CExample}{题目：}}
	已知二次函数 $f(x)=2x^2-4x+1$，求顶点坐标、对称轴方程，并判断 $\Delta$ 符号及零点情况。
	\textbf{\textcolor{CExample}{解题步骤：}}
	\begin{description}[style=nextline, leftmargin=2.8em, labelsep=0.5em, itemsep=0.45em]
		\item[\textbf{第一步（识别系数）：}]
			$a=2,b=-4,c=1$。
			\par\textbf{理由：} 从一般式读出系数。
		\item[\textbf{第二步（计算顶点）：}]
			$h=-\frac{b}{2a}=1$，$k=\frac{4ac-b^2}{4a}=\frac{8-16}{8}=-1$。
			\par\textbf{理由：} 代入顶点坐标公式。
		\item[\textbf{第三步（判别与零点）：}]
			$\Delta=b^2-4ac=16-8=8>0$，零点 $x=\frac{4\pm\sqrt{8}}{4}=1\pm\frac{\sqrt{2}}{2}$。
			\par\textbf{理由：} 判别式为正值表明有两个不同实零点，直接代入求根公式。
	\end{description}
	\textbf{\textcolor{CExample}{关键结论：}}
	\textbf{顶点坐标 $(1,-1)$，对称轴 $x=1$，$\Delta=8>0$，零点 $x=1\pm\frac{\sqrt{2}}{2}$。}

	\medskip
	\textbf{\textcolor{CExample}{例 2（稍变形）}}
	\textbf{\textcolor{CExample}{题目：}}
	已知二次函数的图象与 $x$ 轴交于 $(1,0)$、$(3,0)$，且过点 $(2,-1)$，求解析式。
	\textbf{\textcolor{CExample}{解题步骤：}}
	\begin{description}[style=nextline, leftmargin=2.8em, labelsep=0.5em, itemsep=0.45em]
		\item[\textbf{第一步（设两根式）：}]
			设 $f(x)=a(x-1)(x-3)$。
			\par\textbf{理由：} 与 $x$ 轴交点已知，设两根式。
		\item[\textbf{第二步（代入求 $a$）：}]
			将 $(2,-1)$ 代入得 $-1=a(2-1)(2-3)=a\cdot1\cdot(-1)$，$a=1$。
			\par\textbf{理由：} 代入解得 $a$。
		\item[\textbf{第三步（展开得一般式）：}]
			$f(x)=(x-1)(x-3)=x^2-4x+3$。
			\par\textbf{理由：} 展开即得一般式，也可保留两根式。
	\end{description}
	\textbf{\textcolor{CExample}{关键结论：}}
	\textbf{$f(x)=x^2-4x+3$。}

	\medskip
	\textbf{\textcolor{CExample}{例 3（综合应用）}}
	\textbf{\textcolor{CExample}{题目：}}
	已知二次函数 $f(x)$ 满足 $f(1)=0$，$f(3)=0$，且最大值为 $2$，求解析式。
	\textbf{\textcolor{CExample}{解题步骤：}}
	\begin{description}[style=nextline, leftmargin=2.8em, labelsep=0.5em, itemsep=0.45em]
		\item[\textbf{第一步（设两根式）：}]
			由零点设 $f(x)=a(x-1)(x-3)$。
			\par\textbf{理由：} 根为 $1$ 和 $3$，可设两根式。
		\item[\textbf{第二步（对称轴与最值）：}]
			对称轴 $x=\frac{1+3}{2}=2$，最大值 $f(2)=2$。
			\par\textbf{理由：} 二次函数对称轴为两根中点，且 $a<0$ 才有最大值。
		\item[\textbf{第三步（代入求 $a$）：}]
			$f(2)=a(2-1)(2-3)=a\cdot1\cdot(-1)=-a=2$，得 $a=-2$。
			\par\textbf{理由：} 代入顶点值，解出 $a$。
		\item[\textbf{第四步（展开得一般式）：}]
			\[
				\begin{aligned}
					f(x) &= -2(x-1)(x-3) \\ &= -2(x^2-4x+3) \\ &= -2x^2+8x-6.
			\end{aligned}
			\]
			\par\textbf{理由：} 展开得一般表达式。
	\end{description}
	\textbf{\textcolor{CExample}{关键结论：}}
	\textbf{$f(x)=-2x^2+8x-6$。}
\end{examplebox}
